\section{Lecture 22 -- 10th January 2025}\label{sec: lecture 22}
We state an important corollary of Serre's theorem on cohomology \Cref{thm: Serre cohomology}. 
\begin{corollary}\label{corr: coherence of derived pushforward}
    Let $f:X\to Y$ be a projective morphism of schemes with $Y$ locally Noetherian and $\Fcal$ a coherent sheaf on $X$. Then $f_{*}\Fcal$ and $R^{i}f_{*}\Fcal$ are coherent on $Y$ for all $i\geq 1$.
\end{corollary}
\begin{proof}
    Note that since $f$ is projective, it is in particular quasicompact, so $f_{*}\Fcal$ is a quasicoherent sheaf. The coherence of $f_{*}\Fcal$ is local on target, so without loss of generality take $Y=\spec(A)$ with $A$ Noetherian, in which case we satisfy the hypotheses of \Cref{thm: Serre cohomology}, in which case $f_{*}\Fcal=\widetilde{H^{0}(X,\Fcal)}$ and $R^{i}f_{*}\Fcal=\widetilde{H^{i}(X,\Fcal)}$ -- since both are universal $\delta$-functors hence isomorphic -- are coherent since the $H^{i}$'s are the $\widetilde{(-)}$-ification of finitely generated modules over a Noetherian ring. 
\end{proof}
We have already introduced line bundles -- or locally invertible sheaves -- in \Cref{prop: serre twisting sheaf properties} and \Cref{def: very ample invertible sheaf}. These form a group under $\otimes_{\Ocal_{X}}$ known as the Picard group. 
\begin{definition}[Picard Group]\label{def: Picard group}
    Let $X$ be a scheme. The Picard group $\Pic(X)$ of $X$ is the set of all isomorphism classes of invertible sheaves under the operation $-\otimes_{\Ocal_{X}}-$.
\end{definition}
\begin{remark}
    The inverse of an invertible sheaf $\Lcal$ is the invertible sheaf $\Lcal^{\vee}=\underline{\Hom}_{\Coh(X)}(\Lcal,\Ocal_{X})$. 
\end{remark}
Henceforth, let $X$ be a separated scheme. Recall that a sheaf $\Lcal$ is locally trivial of rank 1 if there exists an open cover $\{U_{i}\}_{i\in I}$ of $X$ on which $\Lcal|_{U_{i}}\cong\Ocal_{U_{i}}$. On double intersections, we have a diagram 
$$% https://q.uiver.app/#q=WzAsMyxbMSwwLCJcXExjYWx8X3tVX3tpan19Il0sWzAsMSwiXFxPY2FsX3tVX3tpan19Il0sWzIsMSwiXFxPY2FsX3tVX3tpan19Il0sWzAsMSwiXFx2YXJwaGlfe2p9fF97VV97aWp9fSIsMl0sWzAsMiwiXFx2YXJwaGlfe2l9fF97VV97aWp9fSJdXQ==
\begin{tikzcd}
	& {\Lcal|_{U_{ij}}} \\
	{\Ocal_{U_{ij}}} && {\Ocal_{U_{ij}}}
	\arrow["{\varphi_{j}|_{U_{ij}}}"', from=1-2, to=2-1]
	\arrow["{\varphi_{i}|_{U_{ij}}}", from=1-2, to=2-3]
\end{tikzcd}$$
where $\varphi_{i}:\Lcal|_{U_{i}}\to\Ocal_{U_{i}},\varphi_{j}:\Lcal|_{U_{j}}\to\Ocal_{U_{j}}$ are the trivializing isomorphisms of $\Lcal$ on the cover. This gives isomorphisms $\varphi_{j}|_{U_{ij}}\circ\varphi_{i}|_{U_{ij}}^{-1}:\Ocal_{U_{ij}}\to\Ocal_{U_{ij}}$ of $\Ocal_{U_{ij}}$-modules. Such an isomorphism is determined by a nowhere vanishing section $\varphi_{ij}(1)\in H^{0}(U_{ij},\Ocal_{U_{ij}}^{\times})$. These satisfy the cocycle conditions: $\varphi_{ii}=1,\varphi_{ij}\circ\varphi_{ji}=1, \varphi_{ij}\circ\varphi_{jk}=\varphi_{ik}$. 

This shows that $\varphi_{ij}\in\check{H}^{1}(X,\Ocal_{X}^{\times})$, and by agreement of \v{C}ech and derived functor cohomoology in degree 1 \Cref{rmk: cech and derived functor in degree 1}, defines an element of $H^{1}(X,\Ocal_{X}^{\times})$. 

Let $\Lcal$ be an invertible sheaf and $s\in H^{0}(X,\Lcal)$. Recall 
$$X_{s}=\{x\in X:s_{x}\notin\mfrak_{x}\Lcal_{x}\}.$$
\begin{proposition}\label{prop: global generation and covering}
    Let $X$ be a scheme and $\Lcal$ an invertible sheaf on $X$. If $\Lcal$ is globally generated by $s_{1},\dots,s_{n}\in H^{0}(X,\Lcal)$ then $\bigcup_{i=1}^{n}X_{s_{i}}=X$. 
\end{proposition}
\begin{proof}
    For $s\in H^{0}(X,\Lcal)$, the morphism $\Ocal_{X}\to\Lcal$ defined by $s$ gives an isomorphism $\Ocal_{X}|_{X_{s}}\to\Lcal|_{X_{s}}$ so $X=\bigcup_{i=1}^{n}X_{s_{i}}$ since $\Lcal$ is globally generated by such sections. 
\end{proof}
In particular, in this situation, we can write $\varphi_{ij}=\frac{s_{j}}{s_{i}}$ on $X_{s_{i}s_{j}}$. In the case of projective space, this gives:
\begin{example}
    Let $A$ be a ring, $X=\PP^{n}_{A}$, and $\Lcal=\Ocal_{X}(1)$. $\Ocal_{X}(1)$ is generated by sections $x_{0},\dots,x_{n}$ and is associated to $\{\frac{x_{j}}{x_{i}}\}_{0\leq i,j\leq n}\in H^{1}(X,\Ocal_{X}^{\times})$. 
\end{example}
This construction holds more generally. 
\begin{proposition}\label{prop: map from Pic to unit torsors}
    Let $X$ be a separated scheme. The map $\Pic(X)\to H^{1}(X,\Ocal_{X}^{\times})$ by $\Lcal\mapsto\{\varphi_{ij}\}_{i,j\in I}$ is an isomorphism of Abelian groups. 
\end{proposition}
\begin{proof}
    We first show that the map is well-defined. Given two open covers on which $\Lcal$ is trivialized, the associated cohomology classes are eventually identified after passing to a sufificiently fine cover in the colimit, hence define the same cohomology class in $H^{1}(X,\Ocal_{X}^{\times})$. 

    The morphism is a bijection of sets since the data of the 1-cocycles $\{\varphi_{ij}\}_{i,j\in I}$ define an invertible sheaf up to isomorphism, showing a bijection between sets. 

    That this is a group isomorphism is obtained by noting that the tensor product of invertible sheaves is mapped to the product of 1-cocycles. 
\end{proof}
\begin{remark}
    There are similar constructions for locally free sheaves of higher rank, but this has bad categorical behavior. 
\end{remark}
\begin{example}
    There is a morphism of groups $\ZZ\to\Pic(\PP^{n}_{A})$ by $d\mapsto\Ocal_{\PP^{n}_{A}}(d)$ which is a group homomorphism and an isomorphism if $A$ is a field. 
\end{example}
This allows us to define maps to projective space over a field. 

Recall that for $f:X\to \PP^{n}_{A}$ an $A$-morphism, there is a map $H^{0}(\PP^{n}_{A},\Ocal_{\PP^{n}_{A}}(1))\to H^{0}(X,f^{*}\Ocal_{\PP^{n}_{A}}(1))$. 
\begin{proposition}\label{prop: maps to Pn}
    \begin{enumerate}[label=(\roman*)]
        \item If $f:X\to\PP^{n}_{A}$ is a morphism of $A$-schemes and $\Lcal\cong f^{*}\Ocal_{\PP^{n}_{A}}(1)$, then the images $s_{0},\dots,s_{n}$ of $x_{0},\dots,x_{n}$ under  $H^{0}(\PP^{n}_{A},\Ocal_{\PP^{n}_{A}}(1))\to H^{0}(X,f^{*}\Ocal_{\PP^{n}_{A}}(1))$ globally generate $\Lcal$. 
        \item Let $X$ be a scheme and $\Lcal$ a line bundle on $X$. If $\Lcal$ is globally generated by $s_{0},\dots,s_{n}$ then there exists a unique morphism $f:X\to\PP^{n}_{A}$ such that $\Lcal=f^{*}\Ocal_{\PP^{n}_{A}}(1)$ and the induced map $H^{0}(\PP^{n}_{A},\Ocal_{\PP^{n}_{A}}(1))\to H^{0}(X,f^{*}\Ocal_{\PP^{n}_{A}}(1))$ takes $x_{0},\dots,x_{n}$ to $s_{0},\dots,s_{n}$. 
    \end{enumerate}
\end{proposition}
\begin{proof}[Proof of (i)]
    We know that on $\PP^{n}_{A}$, the $x_{i}$ genearate $\Ocal_{\PP^{n}_{A}}(1)$ since $\PP^{n}_{A}$ is covered by the distinguished opens $D_{+}(x_{i})$. Now noting that $f^{-1},f^{*}$ are exact and $-\otimes_{f^{-1}\Ocal_{\PP^{n}_{A}}}\Ocal_{X}$ is right exact, we get $\bigoplus_{i=0}^{n}\Ocal_{X}\to\Lcal$ a surjection. 
\end{proof}
\begin{proof}[Proof of (ii)]
    Since $\Lcal$ is globally generated by $s_{0},\dots,s_{n}$, we have that $X=\bigcup_{i=0}^{n}X_{s_{i}}$ by \Cref{prop: global generation and covering}. We define the map locally as maps 
    $$X_{s_{i}}\to D_{+}(x_{i})\cong\spec\left(A\left[\frac{x_{0}}{x_{i}},\dots,\frac{x_{n}}{x_{i}}\right]\right)$$
    which by adjunction with the $\Gamma$-functor is equivalent to specifiying an $A$-algebra map $A[\frac{x_{0}}{x_{i}},\dots,\frac{x_{n}}{x_{i}}]\to\Gamma(X_{s_{i}},\Ocal_{X_{s_{i}}})$ by $\frac{x_{j}}{x_{i}}\mapsto\frac{s_{j}}{s_{i}}$ which glue, giving existence of the morphism. Uniqueness follows from uniqueness of each $A$-algebra morphism. 
\end{proof}
\begin{remark}
    This morphism is a clsoed immersion if and only if each $f^{-1}(D_{+}(x_{i}))$ is affine and the map 
    $$A\left[\frac{x_{0}}{x_{i}},\dots,\frac{x_{n}}{x_{i}}\right]\longrightarrow\Gamma(X_{s_{i}},\Ocal_{X_{s_{i}}})$$
    is surjective. 
\end{remark}
If there is a collection of sections $s_{0},\dots,s_{n}$ that do not globally generate $\Lcal$, then there is still a unique map $\bigcup_{i=0}^{n}X_{s_{i}}\subsetneq X$ to $\PP^{n}_{A}$. This is known as a rational map. 
\begin{example}
    Let $X=\PP^{n+1}_{k}$. The rational map $\PP^{n+1}_{k}\to\PP^{n}_{k}$ by $[x_{0}:\dots:x_{n+1}]\mapsto [x_{0}:\dots:x_{n}]$ is defined by the invertible sheaf generated by sections $x_{0},\dots,x_{n}$ on $\PP^{n+1}_{k}$, which is globally generates $\Ocal_{\PP^{n+1}_{k}}(1)|_{D_{+}(x_{n+1})}$. 
\end{example}
We can now define ample sheaves. 
\begin{definition}[Ample Invertible Sheaf]\label{def: ample invertible sheaf}
    Let $X$ be a scheme and $\Lcal$ an invertible sheaf on $X$. $\Lcal$ is ample if for all coherent sheaves $\Fcal$, $\Fcal\otimes_{\Ocal_{X}}\Lcal^{\otimes n}$ is globally generated for sufificiently large $n$. 
\end{definition}
These are especially easy to understand on affine schemes. 
\begin{example}
    If $X$ is an affine scheme, any invertible sheaf is ample: every coherent sheaf is associated to a finitely presented module. In particular, the structure sheaf of an affine scheme is ample. 
\end{example}